\section{有限覆盖}
\label{sec:metric-finite-covers}

全有界性只回答“能否用有限个球覆盖”，没有回答“需要多少个”。实际问题常常需要这个数量：用一张有限的代表表近似一个集合，精度要求为$\varepsilon$时表要多长？精度提高一倍，表又要变长多少？本节固定一个度量空间$(X,d)$及非空子集$A\subseteq X$，用覆盖数计数所需的代表，用打包数计数能放下多少个彼此分离的点，并说明二者只差一个尺度因子。为了让数量有确定含义，本节一律使用闭球$B[x,\varepsilon]$，并写明球心允许取在哪里。最后用有限维赋范空间的体积比较给出球的覆盖数估计。

\subsection{覆盖数}
\label{sec:metric-covering-number}

先把“用最少的代表覆盖$A$”写成定义。代表即球心，球心是否必须属于$A$会影响数量，所以分两种约定。

\begin{definition}[覆盖数]
\label{def:metric-covering-number}
对$\varepsilon>0$，称
\begin{equation*}
N_X(\varepsilon,A)=\min\Bigl\{m:\text{存在 }x_1,\ldots,x_m\in X\text{ 使 }A\subseteq\bigcup_{j=1}^mB[x_j,\varepsilon]\Bigr\}
\end{equation*}
为$A$在精度$\varepsilon$下的\term{覆盖数}（covering number），不存在这样的有限组时记为$\infty$。要求球心$x_j\in A$时，所得的数记为$N_A(\varepsilon,A)$。满足上述包含关系的点组称为$A$的一个$\varepsilon$网。
\end{definition}

覆盖数随精度单调：$\varepsilon'<\varepsilon$时半径$\varepsilon'$的覆盖也是半径$\varepsilon$的覆盖，所以$N_X(\varepsilon,A)\leq N_X(\varepsilon',A)$；若$A\subseteq A'$，覆盖$A'$的球也覆盖$A$，所以$N_X(\varepsilon,A)\leq N_X(\varepsilon,A')$。两种球心约定之间的差别，只需要把半径调整一倍。

\begin{theorem}[球心范围与覆盖数]
\label{thm:metric-covering-centers}
对任意$\varepsilon>0$，
\begin{equation}
N_X(\varepsilon,A)\leq N_A(\varepsilon,A)\leq N_X(\varepsilon/2,A).
\label{eq:metric-covering-centers}
\end{equation}
\end{theorem}

\begin{proof}
左边的不等式成立，是因为球心在$A$中的覆盖也是球心在$X$中的覆盖。对右边，设$B[x_1,\varepsilon/2],\ldots,B[x_m,\varepsilon/2]$覆盖$A$，并且可以假设每个球都与$A$相交，否则把它删去。在第$j$个球中取$a_j\in A$。对$y\in B[x_j,\varepsilon/2]$，$d(y,a_j)\leq d(y,x_j)+d(x_j,a_j)\leq\varepsilon$，所以$B[x_j,\varepsilon/2]\subseteq B[a_j,\varepsilon]$，球心在$A$中的这$m$个球覆盖$A$。
\end{proof}

由此可知，子集的全有界性不依赖球心约定：$N_A(\varepsilon,A)$对所有$\varepsilon$有限，当且仅当$N_X(\varepsilon,A)$对所有$\varepsilon$有限。数量本身则可能不同。在$\mathbb R$中取$A=\{-1,1\}$，以$0$为中心、半径$1$的闭球覆盖$A$，所以$N_X(1,A)=1$；但球心只能取$\pm1$时，两点相距$2$，必须用两个球，$N_A(1,A)=2$。这时$N_X(1/2,A)=2$，与式\eqref{eq:metric-covering-centers}一致。

再看一个可以完全算出的例子。设$m\geq2$，$\varepsilon=1/(2m)$，$A=[0,1]\subseteq\mathbb R$。以$(2j-1)/(2m)$（$j=1,\ldots,m$）为中心、半径$1/(2m)$的闭区间首尾相接，恰好覆盖$[0,1]$，并且球心都在$A$中，所以$N_A(\varepsilon,A)\leq m$。覆盖数不能更小，这一点下一小节用打包数证明。

\subsection{打包数}
\label{sec:metric-packing-number}

覆盖数的上界只需写出一个覆盖，下界则要说明任何更少的球都不够。打包从相反方向计数：在$A$中放入尽可能多的彼此远离的点，每个小球至多能“照顾”其中一个。

\begin{definition}[分离集与打包数]
\label{def:metric-packing-number}
有限集$S\subseteq A$若满足：对不同的$s,s'\in S$都有$d(s,s')>\varepsilon$，则称$S$为$A$的$\varepsilon$\term{分离集}（separated set）。称
\begin{equation*}
M(\varepsilon,A)=\sup\{\lvert S\rvert:S\text{ 是 }A\text{ 的 }\varepsilon\text{ 分离集}\}
\end{equation*}
为$A$的\term{打包数}（packing number）；存在任意大的有限分离集时，$M(\varepsilon,A)=\infty$。
\end{definition}

\begin{theorem}[覆盖数与打包数]
\label{thm:metric-covering-packing}
对任意$\varepsilon>0$，
\begin{equation}
N_A(\varepsilon,A)\leq M(\varepsilon,A)\leq N_X(\varepsilon/2,A).
\label{eq:metric-covering-packing}
\end{equation}
\end{theorem}

\begin{proof}
左边：若$M(\varepsilon,A)=\infty$，不等式显然。否则取一个元素个数等于$M(\varepsilon,A)$的分离集$S$，它不能再添加$A$中的任何点。于是每个$a\in A$都与某个$s\in S$满足$d(a,s)\leq\varepsilon$，否则$S\cup\{a\}$仍是分离集。所以$\{B[s,\varepsilon]\}_{s\in S}$是球心在$A$中的覆盖，$N_A(\varepsilon,A)\leq\lvert S\rvert$。

右边：设$B[x_1,\varepsilon/2],\ldots,B[x_m,\varepsilon/2]$覆盖$A$，$S$是任意$\varepsilon$分离集。同一个球中的两点距离至多为$\varepsilon$，而$S$中不同点的距离大于$\varepsilon$，所以每个球至多含$S$的一个点；$S\subseteq A$又被这些球覆盖，于是$\lvert S\rvert\leq m$。对$S$取上确界即得。
\end{proof}

左边的论证只用到$S$不能再扩大，所以任何一个“极大”的分离集都是球心在$A$中的$\varepsilon$网，不必是元素最多的那一个。这给出构造网的贪心办法：依次加入与已选点距离都大于$\varepsilon$的点，直到无法再加入为止；若过程能停止，选出的点就是$\varepsilon$网。

现在完成上一小节的例子。设$m\geq2$，点$0,\tfrac1{m-1},\tfrac2{m-1},\ldots,1$共$m$个，两两距离至少为$\tfrac1{m-1}>\tfrac1m$，所以$M(1/m,[0,1])\geq m$。由式\eqref{eq:metric-covering-packing}右边，$N_X(1/(2m),[0,1])\geq M(1/m,[0,1])\geq m$。结合上一小节的覆盖，$N_X(1/(2m),[0,1])=N_A(1/(2m),[0,1])=m$。

把两个定理连起来，
\begin{equation*}
N_X(\varepsilon,A)\leq N_A(\varepsilon,A)\leq M(\varepsilon,A)\leq N_X(\varepsilon/2,A)\leq N_A(\varepsilon/2,A)\leq M(\varepsilon/2,A).
\end{equation*}
三种计数在尺度相差一倍的范围内互相控制。研究数量随$\varepsilon$怎样增长时，可以任选其中最方便的一种。特别地，$A$全有界当且仅当对每个$\varepsilon>0$都有$M(\varepsilon,A)<\infty$。

\subsection{度量熵}
\label{sec:metric-entropy}

覆盖数通常随精度提高而迅速增长，取对数更便于比较增长速度。

\begin{definition}[度量熵]
\label{def:metric-entropy}
在指定的球心约定下，称$H(\varepsilon,A)=\log N_X(\varepsilon,A)$为$A$在精度$\varepsilon$下的\term{度量熵}（metric entropy）；球心限制在$A$中时用$\log N_A(\varepsilon,A)$。
\end{definition}

对数的底只改变一个常数倍。用以$2$为底的对数时，$\lceil\log_2N\rceil$是给$N$个球心编号所需的二进制位数：发送一个编号，接收方取出相应的球心，就能以误差不超过$\varepsilon$复原$A$中的任一点。所以度量熵衡量的是以精度$\varepsilon$描述$A$中元素所需的信息量。

由定义，$A$全有界当且仅当对每个$\varepsilon>0$，$H(\varepsilon,A)<\infty$。这只是有限与否的判断；度量熵进一步记录增长速度。区间$[0,1]$上，$N(1/(2m),[0,1])=m$，所以$H(\varepsilon,[0,1])\approx\log\frac1{2\varepsilon}$。把同样的论证用于$([0,1]^n,\rho_\infty)$：$\rho_\infty$下的闭球是边长$2\varepsilon$的立方体，以$((2j_1-1)/(2m),\ldots,(2j_n-1)/(2m))$为中心的$m^n$个立方体覆盖$[0,1]^n$；坐标取自$\{0,\tfrac1{m-1},\ldots,1\}$的$m^n$个格点两两的$\rho_\infty$距离大于$1/m$，给出相同的下界。于是
\begin{equation*}
N_X\Bigl(\frac1{2m},[0,1]^n\Bigr)=m^n,\qquad H\Bigl(\frac1{2m},[0,1]^n\Bigr)=n\log m.
\end{equation*}
维数作为系数出现在度量熵中：精度提高一倍，所需位数增加约$n$个。对可数无限的离散空间，$\varepsilon<1$时$H=\infty$，它不全有界。

区分“证明有限覆盖存在”与“给出随精度变化的数量估计”很重要。前者只用于判断紧性，例如第~\ref{sec:metric-compactness}节的论证；后者给出可以比较的数值，第~\ref{chap:function-class-covers}章将看到，光滑函数类的度量熵按$1/\varepsilon$的幂增长，远快于有限维情形的$\log(1/\varepsilon)$。

\subsection{有限维球的覆盖}
\label{sec:metric-finite-dim-ball}

上面的立方体例子依赖坐标网格。对一般范数的球，网格不再贴合球的形状，但有限维空间提供了另一种工具：体积。设$(V,\lVert\cdot\rVert)$是$n$维实赋范空间，$B_R=\{\bm x\in V:\lVert\bm x\rVert\leq R\}$。

体积论证需要两条事实。取$V$的一组基，把$V$线性等同于$\mathbb R^n$，用$\mathbb R^n$上的体积测量$V$中的集合；单位球$B_1$是有界闭凸集，并且含有某个欧氏小球（有限维空间上任意两个范数都可以互相控制，见第~\ref{sec:metric-definition}小节引用的卷二章节），所以它的体积$v$满足$0<v<\infty$。其次，平移不改变体积，按因子$r$伸缩使体积乘以$r^n$，所以$\bm x+B_r$的体积为$r^nv$；这是线性换元的Jacobian行列式为$r^n$的结果，可参考\BookChapterRef{chap:region-integrals}{卷二“区域积分”章}的换元公式。

\begin{theorem}[有限维球的覆盖数]
\label{thm:metric-finite-dim-ball-cover}
设$V$是$n$维实赋范空间，$R>0$，$0<\varepsilon\leq R$。则
\begin{equation}
\Bigl(\frac R\varepsilon\Bigr)^n\leq N_V(\varepsilon,B_R)\leq N_{B_R}(\varepsilon,B_R)\leq M(\varepsilon,B_R)\leq\Bigl(1+\frac{2R}\varepsilon\Bigr)^n.
\label{eq:metric-finite-dim-ball-cover}
\end{equation}
\end{theorem}

\begin{proof}
先证最右边的不等式。设$S\subseteq B_R$是$\varepsilon$分离集。对不同的$\bm s,\bm s'\in S$，闭球$\bm s+B_{\varepsilon/2}$与$\bm s'+B_{\varepsilon/2}$不交，否则公共点到两球心的距离都不超过$\varepsilon/2$，与$\lVert\bm s-\bm s'\rVert>\varepsilon$矛盾。这些小球都含于$B_{R+\varepsilon/2}$，因为$\lVert\bm s+\bm y\rVert\leq R+\varepsilon/2$。比较体积，
\begin{equation*}
\lvert S\rvert\Bigl(\frac\varepsilon2\Bigr)^nv\leq\Bigl(R+\frac\varepsilon2\Bigr)^nv,\qquad\text{即}\qquad\lvert S\rvert\leq\Bigl(1+\frac{2R}\varepsilon\Bigr)^n.
\end{equation*}
中间两个不等式就是定理~\ref{thm:metric-covering-centers}与定理~\ref{thm:metric-covering-packing}。

最左边：若$N$个半径$\varepsilon$的闭球覆盖$B_R$，它们的体积之和不小于$B_R$的体积，即$N\varepsilon^nv\geq R^nv$。
\end{proof}

上下界的底数都随$R/\varepsilon$按比例变化，指数都是$n$，所以$H(\varepsilon,B_R)=n\log(R/\varepsilon)+O(n)$。例如$n=1$、$R=1$、$\varepsilon=1/2$时，$[-1,1]$实际只需以$\pm1/2$为中心的两个球，上界$5$、下界$2$都给出正确的数量级，上界只是不够精细。

指数$n$是实维数。复空间$\mathbb C^n$看成实空间时维数是$2n$，相应的上界为$(1+2R/\varepsilon)^{2n}$。证明的关键在于平移与伸缩下体积的变化规律，这来自线性结构和有限维；一般度量空间没有这样的体积。无限维赋范空间中单位球甚至不全有界，第~\ref{chap:normed-spaces}章将说明这一点，那时任何正精度下的覆盖数都是无穷。
